\chapter{Methodology and Model Design}
\label{chap:methodology}

\section{Methodological Overview}

Drug--target interaction (DTI) prediction requires the integration of
heterogeneous biomedical evidence. A drug may be characterized by its chemical
structure, therapeutic associations, side effects, and known target relations,
whereas a protein may be characterized by its amino acid sequence, biological
function, pathway membership, disease associations, and protein--protein
interactions. A biomedical knowledge graph provides a natural representation
for these heterogeneous entities and relations, because drugs, proteins,
diseases, pathways, and other biomedical concepts can be connected through
typed edges.

The original KGE-NFM framework follows a two-stage design
\cite{ye2021unified}. In the first stage, a biomedical knowledge graph is used
to learn low-dimensional representations of drugs, proteins, and other
biomedical entities. In the second stage, the learned drug and protein
representations are combined with structural descriptors and used as input to
a Neural Factorization Machine (NFM) for supervised DTI prediction. The
original framework can be summarized as:

\begin{equation}
    \text{Biomedical KG}
    \rightarrow
    \text{DistMult}
    \rightarrow
    \text{drug/protein KG embeddings}
    \rightarrow
    \text{NFM}
    \rightarrow
    \text{DTI prediction}.
\end{equation}

This report preserves the general KGE + NFM pipeline but modifies the KGE
component. Specifically, the DistMult-based KGE module is replaced by CompGCN,
a composition-based multi-relational graph convolutional network
\cite{vashishth2020composition}. The proposed framework can therefore be
summarized as:

\begin{equation}
    \text{Biomedical KG}
    \rightarrow
    \text{CompGCN}
    \rightarrow
    \text{drug/protein KG embeddings}
    \rightarrow
    \text{NFM}
    \rightarrow
    \text{DTI prediction}.
\end{equation}

The downstream NFM component is retained. The methodological contribution is
located in the KG representation learning stage. This modification is motivated
by the observation that biomedical knowledge graphs are not merely collections
of independent triples. They are directed, multi-relational, and locally
structured graphs. Therefore, the KG encoder should ideally exploit not only
triple-level compatibility, but also relation-aware neighborhood information.

The remainder of this chapter first formalizes the DTI prediction problem and
the two-stage learning objective. It then reviews the original DistMult-based
KGE component, discusses its limitations for directed multi-relational
biomedical graphs, introduces CompGCN as the updated encoder, and finally
presents the complete CompGCN-NFM framework.

\section{Problem Formulation}

\subsection{Drug--Target Interaction Prediction}

Let $d$ denote a drug and $p$ denote a protein target. The DTI prediction task
aims to learn a function that estimates whether $d$ interacts with $p$. Each
candidate pair has a binary label:

\begin{equation}
    y_{dp} \in \{0,1\},
\end{equation}
where $y_{dp}=1$ indicates a known drug--target interaction and $y_{dp}=0$
indicates an unknown or sampled negative interaction. The model outputs an
interaction probability:

\begin{equation}
    \hat{y}_{dp}
    =
    F(d,p),
\end{equation}
where $\hat{y}_{dp}$ denotes the predicted interaction score for the candidate
pair.

\subsection{Biomedical Knowledge Graph}

The biomedical knowledge graph is defined as:

\begin{equation}
    \mathcal{G}
    =
    (\mathcal{E},\mathcal{R},\mathcal{T}),
\end{equation}
where $\mathcal{E}$ is the set of entities, $\mathcal{R}$ is the set of
relation types, and $\mathcal{T}$ is the set of observed triples. Each triple
is represented as:

\begin{equation}
    (h,r,t) \in \mathcal{T},
\end{equation}
where $h,t \in \mathcal{E}$ are the head and tail entities, and
$r \in \mathcal{R}$ is the relation type. In the biomedical setting,
$\mathcal{E}$ may include drugs, proteins, diseases, pathways, side effects,
biological processes, molecular functions, protein domains, and other
biomedical concepts. Relation types may include drug--target interactions,
drug--drug associations, protein--protein interactions, drug--disease
associations, and protein--pathway associations.

The knowledge graph provides auxiliary biomedical context beyond the observed
DTI matrix. This context is especially important because direct DTI
observations are sparse. For example, a protein with few known drug
interactions may still be connected to pathways, biological functions, domains,
and interacting proteins in the KG.

\subsection{Two-Stage Learning Objective}

The framework uses a two-stage learning strategy. Stage 1 learns a KG
representation for each entity:

\begin{equation}
    \mathbf{z}_{e}
    =
    \operatorname{KGE}(e,\mathcal{G}),
    \qquad
    e \in \mathcal{E}.
\end{equation}
Here, $\mathbf{z}_{e}$ is the learned KG embedding of entity $e$. In the
original framework, this representation is obtained using DistMult. In the
proposed framework, it is obtained using CompGCN.

Stage 2 predicts DTI labels using the learned drug and protein embeddings
together with structural descriptors:

\begin{equation}
    \hat{y}_{dp}
    =
    \operatorname{NFM}
    \left(
        \mathbf{z}_{d},
        \mathbf{z}_{p},
        \mathbf{s}_{d},
        \mathbf{s}_{p}
    \right),
\end{equation}
where $\mathbf{z}_{d}$ and $\mathbf{z}_{p}$ are the KG embeddings of drug
$d$ and protein $p$, while $\mathbf{s}_{d}$ and $\mathbf{s}_{p}$ denote drug
and protein structural descriptors. The proposed modification changes how
$\mathbf{z}_{d}$ and $\mathbf{z}_{p}$ are learned, while preserving the NFM
prediction stage.

\section{Knowledge Graph Embedding Background}

\subsection{General KGE Principle}

Knowledge graph embedding is a representation learning approach for
multi-relational data. A KGE model maps entities and relations into a
continuous vector space and defines a scoring function:

\begin{equation}
    f(h,r,t)
\end{equation}
to measure the plausibility of a triple $(h,r,t)$. Training encourages
observed triples to receive higher scores than corrupted triples. For a
positive triple $(h,r,t)$, negative samples can be generated by replacing the
head entity:

\begin{equation}
    (h,r,t)
    \rightarrow
    (h',r,t),
\end{equation}
or by replacing the tail entity:

\begin{equation}
    (h,r,t)
    \rightarrow
    (h,r,t').
\end{equation}

The KGE loss can be written as a margin-based ranking objective:

\begin{equation}
    \mathcal{L}_{KGE}
    =
    \sum_{(h,r,t)\in\mathcal{T}^{+}}
    \sum_{(h',r,t')\in\mathcal{T}^{-}}
    \max
    \left(
        0,
        \gamma
        +
        f(h',r,t')
        -
        f(h,r,t)
    \right),
    \label{eq:kge_margin_loss}
\end{equation}
where $\mathcal{T}^{+}$ denotes observed triples, $\mathcal{T}^{-}$ denotes
sampled corrupted triples, and $\gamma$ is the margin. The corresponding
parameter update is:

\begin{equation}
    \Theta_{KGE}^{(q+1)}
    =
    \Theta_{KGE}^{(q)}
    -
    \eta
    \nabla_{\Theta_{KGE}}
    \mathcal{L}_{KGE},
    \label{eq:kge_update}
\end{equation}
where $\Theta_{KGE}$ denotes KGE parameters and $\eta$ is the learning rate.

\subsection{Role of KGE in DTI Prediction}

In DTI prediction, KGE is used to summarize the biomedical context of drugs
and proteins. The embedding of a protein may encode information about related
pathways, biological functions, diseases, domains, and interacting proteins.
The embedding of a drug may encode information about related targets,
diseases, side effects, therapeutic classes, and drug--drug relations.

These KG-derived representations complement structural features such as
chemical fingerprints and protein sequence descriptors. Within the original
KGE-NFM framework, this representation learning stage is implemented using
DistMult, a bilinear knowledge graph embedding model known for its simplicity
and computational efficiency \cite{yang2015embedding}.

\section{Original KGE Component: DistMult}

\subsection{DistMult Representation Assumption}

DistMult assigns a vector representation to each entity and relation. For a
triple $(h,r,t)$, the head entity, relation, and tail entity are represented
as:

\begin{equation}
    \mathbf{h}, \mathbf{r}, \mathbf{t}
    \in
    \mathbb{R}^{k}.
\end{equation}

DistMult can be viewed as a simplified bilinear model. A full bilinear model
uses a relation-specific matrix $\mathbf{M}_{r}$:

\begin{equation}
    f_{\text{bilinear}}(h,r,t)
    =
    \mathbf{h}^{\top}
    \mathbf{M}_{r}
    \mathbf{t}.
\end{equation}

DistMult restricts the relation matrix to be diagonal:

\begin{equation}
    \mathbf{M}_{r}
    =
    \operatorname{diag}(\mathbf{r}).
\end{equation}
This assumption greatly reduces the number of trainable parameters for each
relation.

\subsection{DistMult Scoring Function}

The DistMult scoring function is:

\begin{equation}
    f_{\text{DM}}(h,r,t)
    =
    \mathbf{h}^{\top}
    \operatorname{diag}(\mathbf{r})
    \mathbf{t}.
    \label{eq:distmult_score_matrix}
\end{equation}
Equivalently, it can be written as a dimension-wise multiplicative interaction:

\begin{equation}
    f_{\text{DM}}(h,r,t)
    =
    \sum_{i=1}^{k}
    h_i r_i t_i.
    \label{eq:distmult_score_sum}
\end{equation}

Equations~\ref{eq:distmult_score_matrix} and
\ref{eq:distmult_score_sum} show that the score is computed by matching the
same embedding dimensions of the head and tail entities under relation-specific
weights. This design avoids using a full $k \times k$ relation matrix and
therefore reduces computational and parameter cost.

\subsection{Rationale for Using DistMult in the Original Framework}

DistMult is a reasonable KGE choice for the original KGE-NFM framework because
it is simple, scalable, and capable of producing compact embeddings for
entities and relations. Its diagonal bilinear form is suitable for relatively
large biomedical KGs, where learning a full matrix for every relation would be
expensive. The output entity embeddings can also be directly used as input
features for the downstream NFM module.

In the original framework, DistMult is not used as the final DTI classifier.
Rather, it is used to extract KG-based drug and protein representations:

\begin{equation}
    \mathbf{z}_{d}
    =
    \operatorname{DistMult}(d),
    \qquad
    \mathbf{z}_{p}
    =
    \operatorname{DistMult}(p).
\end{equation}

Nevertheless, the same assumptions that make DistMult efficient also restrict
its expressiveness when the underlying graph contains directed and
semantically diverse biomedical relations.

\section{Limitations of DistMult and Motivation for the Update}

\subsection{Symmetric Scoring and Directionality Limitation}

The DistMult score is symmetric with respect to the head and tail entities:

\begin{equation}
    f_{\text{DM}}(h,r,t)
    =
    f_{\text{DM}}(t,r,h).
    \label{eq:distmult_symmetry}
\end{equation}

This property follows from the commutativity of element-wise multiplication in
Equation~\ref{eq:distmult_score_sum}. Although this symmetry can be acceptable
for some relations, it creates a structural mismatch for biomedical KGs. Many
biomedical relations are directional or asymmetric. For example,
\textit{drug targets protein}, \textit{drug treats disease}, and
\textit{protein participates in pathway} are not semantically equivalent to
their reversed forms. A scoring function that is symmetric for a fixed relation
cannot explicitly distinguish these directional roles.

\subsection{Lack of Explicit Neighborhood Aggregation}

DistMult learns embeddings through triple-level scoring, but it does not
explicitly aggregate information from the local neighborhood of an entity. This
is important because the biomedical meaning of a drug or protein is often
contextual. A protein may be characterized by its related pathways, domains,
diseases, and interacting proteins. A drug may be characterized by its targets,
side effects, therapeutic classes, and related diseases.

For DTI prediction, it is therefore desirable for the KG encoder to update
entity representations using relation-aware neighborhood information rather
than relying only on independent triple scoring. This limitation motivates a
graph encoder that treats the KG as a structured multi-relational graph rather
than only as a set of independent triples.

\subsection{Limited Use of Relation Embeddings}

In DistMult, relation embeddings mainly act as parameters inside the scoring
function. They determine the dimension-wise compatibility between a head entity
and a tail entity, but they do not explicitly define how messages from
neighboring entities should be transformed or propagated.

This is restrictive for multi-relational biomedical graphs because different
relation types should contribute different kinds of evidence. A drug connected
to a disease relation and a drug connected to a target relation should not
receive the same type of relational signal. Similarly, protein--pathway and
protein--protein relations should affect protein representations differently.

\subsection{Requirements for an Improved KGE Encoder}

The preceding limitations suggest that an improved KGE encoder should preserve
relation direction, distinguish relation types, jointly represent entities and
relations, incorporate local graph neighborhood structure, and produce
embeddings compatible with the downstream NFM module.

These considerations motivate the replacement of DistMult with a relational
graph encoder. In this work, CompGCN is adopted because it directly
incorporates relation type and edge direction into the message-passing process
while jointly updating entity and relation representations.

\section{Updated KGE Component: CompGCN}

\subsection{Overview of CompGCN}

CompGCN is a multi-relational graph convolutional framework designed for
knowledge graphs \cite{vashishth2020composition}. In contrast to DistMult,
CompGCN does not learn entity embeddings only through isolated triple scores.
Instead, it updates each entity representation by aggregating messages from its
relational neighbors. The message from a neighbor is conditioned on both the
neighboring entity and the relation connecting that neighbor to the target
node.

This design is suitable for the present framework because biomedical KGs
contain relation labels and edge directions that carry biological meaning.
CompGCN jointly models entities and relations and uses a composition function
to combine neighboring entity embeddings with relation embeddings before
message passing.

\subsection{Graph Augmentation with Directional Edges}

For every original edge:

\begin{equation}
    (u,r,v),
\end{equation}
CompGCN introduces an inverse edge:

\begin{equation}
    (v,r^{-1},u),
\end{equation}
and a self-loop relation:

\begin{equation}
    (u,\top,u).
\end{equation}

Original edges preserve the observed direction of the relation. Inverse edges
allow information to flow in the reverse direction. Self-loops preserve the
node's own information during message passing. This augmentation directly
addresses the directionality limitation of DistMult by allowing the model to
distinguish original, inverse, and self-loop messages.

\subsection{Entity--Relation Composition}

The message sent from a neighboring entity is relation-dependent. For a
neighbor $u$ connected through relation $r$, CompGCN first computes:

\begin{equation}
    \phi
    \left(
        \mathbf{h}_{u}^{k},
        \mathbf{h}_{r}^{k}
    \right),
\end{equation}
where $\mathbf{h}_{u}^{k}$ is the embedding of entity $u$ at layer $k$,
$\mathbf{h}_{r}^{k}$ is the embedding of relation $r$ at layer $k$, and
$\phi$ is an entity--relation composition operator.

Common composition choices include subtraction:

\begin{equation}
    \phi(\mathbf{h}_{u},\mathbf{h}_{r})
    =
    \mathbf{h}_{u}
    -
    \mathbf{h}_{r},
\end{equation}
element-wise multiplication:

\begin{equation}
    \phi(\mathbf{h}_{u},\mathbf{h}_{r})
    =
    \mathbf{h}_{u}
    \odot
    \mathbf{h}_{r},
\end{equation}
and circular correlation:

\begin{equation}
    \phi(\mathbf{h}_{u},\mathbf{h}_{r})
    =
    \mathbf{h}_{u}
    \star
    \mathbf{h}_{r}.
\end{equation}

This composition operation allows the same neighboring entity to contribute
different messages under different relation types. In this project, the main
CompGCN configuration uses element-wise multiplication, which preserves the
multiplicative relational interaction used by DistMult while embedding it
inside a graph convolutional update.

\subsection{CompGCN Node Update}

The CompGCN node update is:

\begin{equation}
    \mathbf{h}_{v}^{k+1}
    =
    \sigma
    \left(
        \sum_{(u,r)\in \mathcal{N}(v)}
        \mathbf{W}_{\lambda(r)}^{k}
        \phi
        \left(
            \mathbf{h}_{u}^{k},
            \mathbf{h}_{r}^{k}
        \right)
    \right).
    \label{eq:compgcn_node_update}
\end{equation}

In Equation~\ref{eq:compgcn_node_update}, $\mathcal{N}(v)$ denotes the set of
relational neighbors of node $v$, $\mathbf{h}_{u}^{k}$ is the embedding of
neighbor $u$ at layer $k$, $\mathbf{h}_{r}^{k}$ is the embedding of relation
$r$, and $\phi$ composes the neighbor and relation embeddings.
$\mathbf{W}_{\lambda(r)}^{k}$ is a direction-specific transformation matrix,
and $\sigma$ is a nonlinear activation function. The function $\lambda(r)$
identifies whether the edge is original, inverse, or self-loop. Therefore, the
model can transform messages differently depending on edge direction.

Compared with DistMult, this update introduces an explicit mechanism for
incorporating relation-aware neighborhood structure into entity
representations.

\subsection{CompGCN Relation Update}

CompGCN also updates relation embeddings across layers:

\begin{equation}
    \mathbf{h}_{r}^{k+1}
    =
    \mathbf{W}_{rel}^{k}
    \mathbf{h}_{r}^{k}.
    \label{eq:compgcn_relation_update}
\end{equation}

Equation~\ref{eq:compgcn_relation_update} keeps relation representations
aligned with the evolving entity embedding space. Thus, relation embeddings are
not static labels; they remain part of the representation learning process.

\subsection{Link Prediction Decoder}

CompGCN serves as an encoder. After message passing, it produces updated
entity and relation embeddings. A decoder is then used to score triples:

\begin{equation}
    f(h,r,t)
    =
    \operatorname{Dec}
    \left(
        \mathbf{h}_{h}^{L},
        \mathbf{h}_{r}^{L},
        \mathbf{h}_{t}^{L}
    \right).
\end{equation}

To maintain a controlled comparison with the original framework, the decoder
can be written in a DistMult-style form:

\begin{equation}
    f(h,r,t)
    =
    \mathbf{h}^{\top}
    \operatorname{diag}(\mathbf{r})
    \mathbf{t}.
    \label{eq:compgcn_distmult_decoder}
\end{equation}

The original KGE module consists of trainable entity and relation embeddings
with DistMult scoring. The updated KGE module uses a CompGCN encoder to
generate relation-aware entity and relation embeddings before triple scoring.
The key improvement is therefore in how the embeddings are generated, not in
replacing the downstream DTI classifier.

\section{Construction of DTI Pair Features}

\subsection{Extracting Drug and Protein KG Embeddings}

After the KGE stage, each drug $d$ and protein $p$ has a learned embedding:

\begin{equation}
    \mathbf{z}_{d},
    \mathbf{z}_{p}.
\end{equation}

For a candidate DTI pair, the KG-based feature is constructed as:

\begin{equation}
    \mathbf{x}_{KG}(d,p)
    =
    [
    \mathbf{z}_{d}
    \mathbin{\Vert}
    \mathbf{z}_{p}
    ],
    \label{eq:kg_pair_feature}
\end{equation}
where $\mathbin{\Vert}$ denotes concatenation. In the original framework,
these embeddings are obtained from DistMult. In the proposed framework, they
are generated by CompGCN.

\subsection{Incorporating Structural Descriptors}

KG embeddings capture relational biomedical context, but DTI prediction may
also benefit from intrinsic drug and protein information. Let
$\mathbf{s}_{d}$ denote the drug structural descriptor, such as Morgan
fingerprints, and let $\mathbf{s}_{p}$ denote the protein descriptor, such as
CTD or other sequence-derived descriptors. The final input representation is:

\begin{equation}
    \mathbf{x}_{dp}
    =
    [
    \mathbf{z}_{d}
    \mathbin{\Vert}
    \mathbf{z}_{p}
    \mathbin{\Vert}
    \mathbf{s}_{d}
    \mathbin{\Vert}
    \mathbf{s}_{p}
    ].
    \label{eq:dti_pair_feature}
\end{equation}

Equation~\ref{eq:dti_pair_feature} combines heterogeneous KG information with
structural information. The KG embeddings summarize graph context, while the
drug and protein descriptors provide direct molecular and sequence-level
evidence.

\section{Neural Factorization Machine for DTI Prediction}

\subsection{Motivation for Using NFM}

The final feature vector contains multiple heterogeneous feature groups. Their
interactions are important for predicting DTI. For example, the drug KG
embedding may interact with the protein KG embedding, the drug fingerprint may
interact with the protein descriptor, and graph-derived drug context may
interact with sequence-derived protein information.

NFM is used because it explicitly models second-order feature interactions
through Bi-Interaction pooling and further captures nonlinear higher-order
patterns through neural layers \cite{he2017neural}. This design preserves the
original KGE-NFM philosophy: KG embeddings provide biomedical context, while
NFM learns how this context interacts with structural features.

\subsection{NFM Prediction Function}

Given an input feature vector $\mathbf{x}$, the NFM prediction function is:

\begin{equation}
    \hat{y}(\mathbf{x})
    =
    w_0
    +
    \sum_{i=1}^{n}
    w_i x_i
    +
    f(\mathbf{x}),
    \label{eq:nfm_prediction}
\end{equation}
where $w_0$ is the global bias, $w_i$ models the first-order contribution of
feature $i$, and $f(\mathbf{x})$ represents the nonlinear interaction
component. The first two terms preserve the linear part of factorization
machines, while $f(\mathbf{x})$ models feature interactions.

\subsection{Bi-Interaction Pooling}

Let $\mathbf{v}_i$ denote the embedding vector of feature $i$. NFM constructs a
set of scaled feature embeddings and applies Bi-Interaction pooling:

\begin{equation}
    f_{BI}(\mathcal{V}_{x})
    =
    \sum_{i=1}^{n}
    \sum_{j=i+1}^{n}
    x_i\mathbf{v}_i
    \odot
    x_j\mathbf{v}_j.
    \label{eq:nfm_bi_pooling}
\end{equation}

The operator $\odot$ denotes element-wise multiplication. This layer summarizes
pairwise interactions among features into a single vector, which is then
passed into a multi-layer neural network.

\subsection{Hidden Layers and Final Prediction}

The hidden layer transformation is:

\begin{equation}
    \mathbf{z}_{1}
    =
    \sigma
    \left(
        \mathbf{W}_{1}
        f_{BI}
        +
        \mathbf{b}_{1}
    \right),
\end{equation}
and for deeper layers:

\begin{equation}
    \mathbf{z}_{l}
    =
    \sigma
    \left(
        \mathbf{W}_{l}
        \mathbf{z}_{l-1}
        +
        \mathbf{b}_{l}
    \right).
\end{equation}

The final binary prediction is:

\begin{equation}
    \hat{y}_{dp}
    =
    \operatorname{sigmoid}
    \left(
        \mathbf{p}^{\top}
        \mathbf{z}_{L}
    \right).
    \label{eq:nfm_final_prediction}
\end{equation}

Equation~\ref{eq:nfm_final_prediction} gives the probability that the
candidate drug--protein pair interacts. During supervised training, the NFM
parameters are optimized with binary cross-entropy:

\begin{equation}
    \mathcal{L}_{DTI}
    =
    -
    \frac{1}{N}
    \sum_{i=1}^{N}
    \left[
        y_i \log(\hat{y}_i)
        +
        (1-y_i)
        \log(1-\hat{y}_i)
    \right].
    \label{eq:dti_bce_loss}
\end{equation}

\section{Final Proposed CompGCN-NFM Framework}

\subsection{Complete Pipeline}

The final method is a stage-wise framework. First, a biomedical knowledge graph
is constructed from DTI and supporting biomedical relations. The graph is then
augmented with inverse and self-loop edges so that direction-aware message
passing can be performed. CompGCN is trained to learn relation-aware
representations of drugs, proteins, and other biomedical entities. After
training, the drug and protein embeddings are extracted from the CompGCN
encoder. These embeddings are concatenated with drug and protein structural
descriptors to form the DTI pair representation. The final feature vector is
then passed into NFM, which learns feature interactions and predicts the binary
DTI label.

The procedure can be summarized as:

\begin{enumerate}
    \item Construct a biomedical knowledge graph from DTI and supporting
    biomedical relations.
    \item Augment the graph with inverse and self-loop edges.
    \item Train CompGCN to learn relation-aware drug and protein embeddings.
    \item Extract drug and protein embeddings from the trained CompGCN encoder.
    \item Concatenate these embeddings with drug and protein structural
    descriptors.
    \item Feed the final DTI pair feature vector into NFM.
    \item Train the NFM model for binary DTI prediction.
\end{enumerate}

\subsection{Difference from the Original KGE-NFM Framework}

Table~\ref{tab:distmult_compgcn_update} summarizes the methodological
difference between the original KGE-NFM framework and the proposed
CompGCN-NFM framework.

\begin{table}[H]
\centering
\caption{Comparison between original KGE-NFM and proposed CompGCN-NFM}
\label{tab:distmult_compgcn_update}
\small
\begin{tabularx}{\textwidth}{lXX}
\toprule
\textbf{Component} & \textbf{Original KGE-NFM} & \textbf{Proposed CompGCN-NFM} \\
\midrule
KG representation model
& DistMult
& CompGCN \\
Entity representation learning
& Triple-level scoring
& Relation-aware message passing \\
Relation direction
& Limited by symmetric scoring
& Original and inverse relations are explicitly modeled \\
Neighborhood structure
& Not explicitly aggregated
& Aggregated through graph convolution \\
Relation embedding role
& Used mainly in triple scoring
& Used in entity--relation message composition \\
Downstream DTI predictor
& NFM
& NFM \\
Structural descriptors
& Used
& Used \\
\bottomrule
\end{tabularx}
\end{table}

The proposed framework does not discard the original KGE-NFM design. Instead,
it retains the NFM-based feature interaction module and modifies the KG
representation stage. This allows the framework to benefit from the original
multimodal integration strategy while using a more expressive encoder for
directed multi-relational biomedical graphs.

\subsection{Expected Benefit in Cold-Start Scenarios}

Cold-start DTI prediction is difficult because some drugs or proteins have no
observed DTI links in the training set. However, these entities may still have
other biomedical relations in the KG. Since CompGCN aggregates information
from relational neighborhoods, it may produce more informative embeddings for
such entities than a purely triple-scoring KGE model.

The proposed update is expected to be particularly useful in cold-start
scenarios, where direct DTI evidence is sparse but auxiliary KG relations
remain available. This expectation is methodological rather than a substitute
for empirical validation; the experimental chapter evaluates whether the
proposed representation update improves predictive performance.

\section{Model Configuration and Implementation Details}

The implementation uses PyKEEN for KGE training and a local PyTorch
implementation of NFM for supervised DTI prediction. Table~\ref{tab:model_configuration}
summarizes the main configuration used in this report.

\begin{table}[H]
\centering
\caption{Model configuration and implementation details}
\label{tab:model_configuration}
\small
\begin{tabularx}{\textwidth}{lX}
\toprule
\textbf{Item} & \textbf{Configuration} \\
\midrule
KG entity types
& Drugs, proteins, genes, pathways, functions, domains, families, diseases,
side effects, and related biomedical concepts \\
KG relation types
& Drug--target, drug--drug, protein--protein, drug--disease,
protein--pathway, annotation, domain, family, and functional relations \\
Yamanishi08 KG size
& 25,487 entities, 48 relations, and 95,672 triples \\
BioKG size
& 101,847 entities, 13 relations, and 2,017,795 triples \\
KGE baseline
& DistMult \\
Proposed KGE encoder
& CompGCN \\
Decoder
& DistMult-style scoring for controlled KGE comparison \\
Embedding dimension
& 400 \\
CompGCN layers
& 2--3 depending on experiment setting \\
Composition operator
& Element-wise multiplication, \texttt{mult} \\
Negative sampling
& Basic negative sampling during KGE training \\
KGE optimizer and loss
& Adam optimizer with margin ranking loss \\
Drug descriptor
& Morgan fingerprint \\
Protein descriptor
& CTD sequence descriptor \\
NFM hidden layer sizes
& 128, 128 \\
Activation function
& ReLU in NFM hidden layers \\
Dropout
& 0.0--0.2 depending on experiment setting \\
NFM optimizer and loss
& Adam optimizer with binary cross-entropy objective \\
Learning rate
& 0.001 for KGE and NFM unless otherwise specified \\
Batch size
& KGE batch size 1024; NFM batch size up to 20000 \\
Evaluation metrics
& ROC-AUC and PR-AUC \\
\bottomrule
\end{tabularx}
\end{table}

The evaluation settings include warm start, cold start for drugs, and cold
start for proteins. In the warm-start setting, drugs and proteins in the test
set also appear in the training data. In the drug cold-start setting, test
drugs do not appear in the positive training DTI pairs. In the protein
cold-start setting, test proteins do not appear in the positive training DTI
pairs. These settings are used to evaluate whether the proposed KG
representation update supports generalization beyond observed DTI links.
